\chapter*{Практическая часть}

\section*{Задание 1}
Запустить среду Visual Prolog 5.2. Настроить утилиту TestGoal.
Запустить тестовую программу, проанализировать реакцию системы и множество ответов.
Разработать свою программу - «Телефонный справочник».
Протестировать работу программы.

При установке утилиты TestGoal добиться того, чтобы в результате выполнения
тестовой программы возвращалось несколько вариантов ответов.

\section*{Задание 2}
Разработать свою программу - «Телефонный справочник и автомобили». Абоненты
могут иметь несколько телефонов. Протестировать работу программы, используя разные вопросы.

\begin{itemize}
    \item «Телефонный справочник»: Фамилия, Номер, Адрес - структура (Город, Улица, Nºдома, N°кв)
    \item «Автомобили»: Фамилия владельца, Марка, Цвет, Стоимость, Номер.
\end{itemize}

Владелец может иметь несколько телефонов, автомобилей (Факты). В разных городах есть однофамильцы, в одном городе - фамилия уникальна.
Используя коньюнктивное правило и простой вопрос, обеспечить возможность
поиска:

По Марке и Цвету автомобиля найти Фамилию, Город, Телефон.

На листинге~\ref{lst:brt} представлен код программы:
%\FloatBarrier
%\begin{lstinputlisting}[style={lsp}]{src/main2.pro}
%\end{lstinputlisting}
%\FloatBarrier

\begin{center}
    \captionsetup{justification=raggedright,singlelinecheck=off}
    \begin{lstlisting}[label=lst:brt,caption=Код программы]
domains
    name = symbol.
    city = symbol.
    street = symbol.
    car_brand = symbol.
    color = symbol.
    phone = string.
    address = addressStruct(city, street, integer, integer).

predicates
    phone_number(name, address, phone).
    car_info(name, car_brand, color, integer, symbol).

    find_owner_by_car(car_brand, color, name, city, phone).

clauses

    phone_number(ivanov,
        addressStruct(moscow, bauman, 10, 5),
            "8-926-987-65-43").
    phone_number(popov,
        addressStruct(tula, pervomayskaya, 22, 15),
            "8-903-555-22-11").
    phone_number(ivanov,
        addressStruct(volgograd, davidkovo, 3, 2),
            "8-999-777-88-66").

    car_info(ivanov, toyota, blue, 1000000, a123aa).
    car_info(popov, bmw, black, 2000000, b456bb).
    car_info(ivanov, lada, red, 500000, c789cc).

    find_owner_by_car(Brand, Color, Name, City, Phone) :-
        car_info(Name, Brand, Color, _, _),
        phone_number(Name, addressStruct(City, _, _, _),
                                                Phone).

goal
	find_owner_by_car(toyota, blue, Name, City, Phone).
	%find_owner_by_car(tesla, white, Name, City, Phone).
	%find_owner_by_car(lada, red, Name, City, Phone).
	%car_info(Name, Brand, Color, _, _).
    \end{lstlisting}
\end{center}

\chapter*{Тестирование}
\begin{center}
    \captionsetup{justification=raggedright,singlelinecheck=off}
    \begin{lstlisting}[label=lst:brute_fRRRorce,caption=Тестирование]
find_owner_by_car(toyota, blue, Name, City, Phone).
Name=ivanov, City=moscow, Phone=8-926-987-65-43
Name=ivanov, City=volgograd, Phone=8-999-777-88-66
2 Solutions

find_owner_by_car(tesla, white, Name, City, Phone).
No Solution

find_owner_by_car(bmw, black, Name, City, Phone).
Name=popov, City=tula, Phone=8-903-555-22-11
1 Solution
    \end{lstlisting}
\end{center}

\chapter*{Теоретические вопросы}
Основным элементом языка Prolog является терм: константа, переменная или составной терм.
В некоторых случаях, можно сказать, что составной терм является предикатом.

Программа на Prolog не является последовательностью действий, - она представляет
собой набор фактов и правил, которые формируют базу знаний о предметной области.
Факты представляют собой составные термы, с помощью которых фиксируется наличие
истинностных отношений между объектами предметной области — аргументами терма.

Утверждения программы — это предикаты. Предикаты могут не содержать переменных
(основные) или содержать переменные (не основные). В процессе выполнения программы —
система пытается найти, используя базу знаний , такие значения переменных, при которых на
поставленный вопрос можно дать ответ «Да».

Программа <<Телефонный справочник>> на Visual Prolog --- база знаний, которая хранит информацию о людях,
их телефонных номерах, адресах и автомобилях, позволяет находить людей по характеристикам автомобиля (модель, цвет) и
получить связанные с ним данные (фамилия, номер телефона, город) и вопрос, который воспринимается как цель, которую
нужно достичь, используя знания из базы.

Структура программы:
\begin{itemize}
    \item Домены. Домены определяют типы данных: car brand, color, street, city, name, phone, address;
    \item Предикаты: phone number(name, address, phone) - хранит информацию о телефонах и адресах; car info(name, car brand, color, integer, symbol) - хранит данные об автомобилях;
    find owner by car(Brand, Color, Name, City, Phone) - связывает автомобиль с городом и телефоном по фамилии
    \item Факты. Список фактов для phone number и car info - хранятся конкретные записи;
    \item Правила. find owner by car(Brand, Color, Name, City, Phone) - Объединяет данные из car info и phone number через общую переменную.
\end{itemize}
